\documentclass[a4paper,12pt]{article}
\usepackage[english]{babel}
\usepackage[utf8]{inputenc}
\usepackage[T1]{fontenc}
\usepackage{enumitem}
\usepackage{amsmath,amssymb}
\usepackage{geometry}
\geometry{left=25mm,right=25mm,top=25mm,bottom=25mm}
\usepackage{parskip}

\usepackage{tikz}
\usetikzlibrary{circuits.ee.IEC, positioning}

\newif\ifshowsolutions
\showsolutionstrue  % comment to hide solutions

\begin{document}

Kreuze jeweils eine Antwort pro Frage an.

\vspace{8pt}

\begin{enumerate}[label=\textbf{\arabic*.}, leftmargin=*, itemsep=8pt]

  \item Welche Größe ist maßgeblich für die thermische Belastung eines stromdurchflossenen Leiters?
    \begin{itemize}[label=$\square$]
      \item[A)] die elektrische Feldstärke
      \item[B)] die elektrische Leistung
      \item[C)] die elektrische Stromdichte
      \item[D)] der elektrische Widerstand
    \end{itemize}

\ifshowsolutions
Lösung: C
\fi

  \item Ein idealer ohmscher Widerstand ist dadurch gekennzeichnet, dass
    \begin{itemize}[label=$\square$]
      \item[A)] Strom und Spannung proportional zueinander sind.
      \item[B)] eine Phasenverschiebung zwischen Strom und Spannung auftritt.
      \item[C)] der Widerstandswert zeitabhängig ist.
      \item[D)] elektrische Energie periodisch gespeichert und abgegeben wird.
    \end{itemize}

\ifshowsolutions
Lösung: A
\fi

  \item Welche Aussage beschreibt die $I$-$U$-Kennlinie einer idealen Stromquelle?
    \begin{itemize}[label=$\square$]
      \item[A)] Die Spannung ist proportional zum Strom.
      \item[B)] Der Strom ist unabhängig von der Klemmenspannung.
      \item[C)] Die Kennlinie verläuft durch den Ursprung.
      \item[D)] Die Kennlinie ist frequenzabhängig.
    \end{itemize}

\ifshowsolutions
Lösung: B
\fi

  \item Welche Aussage über elektrische Leistung und Energie ist richtig?
    \begin{itemize}[label=$\square$]
      \item[A)] Leistung und Energie sind unabhängig voneinander.
      \item[B)] Energie ist die zeitliche Ableitung der Leistung.
      \item[C)] Bei konstanter Leistung wächst die Energie linear mit der Zeit.
      \item[D)] Leistung und Energie unterscheiden sich nur durch einen konstanten Faktor.
    \end{itemize}

\ifshowsolutions
Lösung: C
\fi

  \item Die Maschenregel ist eine Konsequenz von:
    \begin{itemize}[label=$\square$]
      \item[A)] Ladungserhaltung
      \item[B)] Impulserhaltung
      \item[C)] Temperaturerhaltung
      \item[D)] Energieerhaltung
    \end{itemize}

\ifshowsolutions
Lösung: D
\fi

  \item Eine Y–$\Delta$-Transformation ist in einem Gleichstromnetz zulässig, wenn
    \begin{itemize}[label=$\square$]
      \item[A)] alle Widerstände denselben Wert besitzen.
      \item[B)] das Netzwerk nur von einer Spannungsquelle gespeist wird.
      \item[C)] keine Leistung im Netzwerk umgesetzt wird.
      \item[D)] die Strom-Spannungs-Beziehungen an den äußeren Klemmen unverändert bleiben.
    \end{itemize}

\ifshowsolutions
Lösung: D
\fi

  \item Ein Voltmeter besitzt idealerweise:
    \begin{itemize}[label=$\square$]
      \item[A)] kleinen Innenwiderstand
      \item[B)] großen Innenwiderstand
      \item[C)] keinen Innenwiderstand
      \item[D)] frequenzabhängigen Widerstand
    \end{itemize}

\ifshowsolutions
Lösung: B
\fi

  \item Welche Aussage beschreibt den Wirkungsgrad eines elektrischen Systems korrekt?
    \begin{itemize}[label=$\square$]
      \item[A)] Er ist das Verhältnis von abgegebener zu aufgenommener Leistung.
      \item[B)] Er ist das Verhältnis von Blindleistung zur Wirkleistung.
      \item[C)] Er ist unabhängig von den elektrischen Verlusten im System.
      \item[D)] Er kann im Wechselstrombetrieb größer als eins sein.
    \end{itemize}

\ifshowsolutions
Lösung: A
\fi

  \item Welche Aussage über elektrische und magnetische Feldlinien ist korrekt?
    \begin{itemize}[label=$\square$]
      \item[A)] Elektrische Feldlinien sind stets geschlossen.
      \item[B)] Magnetische Feldlinien sind stets geschlossen.
      \item[C)] Magnetische Feldlinien enden an magnetischen Ladungen.
      \item[D)] Elektrische Feldlinien existieren nur im Leiter.
    \end{itemize}

\ifshowsolutions
Lösung: B
\fi

  \item Welche der folgenden Größen ist ein direkt messbares physikalisches Feld in der makroskopischen Elektrodynamik?
    \begin{itemize}[label=$\square$]
      \item[A)] elektrische Verschiebung \(\mathbf{D}\)
      \item[B)] magnetische Feldstärke \(\mathbf{H}\)
      \item[C)] magnetische Flussdichte \(\mathbf{B}\)
      \item[D)] elektrische Polarisation \(\mathbf{P}\)
    \end{itemize}

\ifshowsolutions
Lösung: C
\fi

  \item Ein zeitlich veränderliches magnetisches Feld erzeugt gemäß Faradays Gesetz:
    \begin{itemize}[label=$\square$]
      \item[A)] eine induzierte Spannung
      \item[B)] einen konstanten Strom
      \item[C)] Wärme
      \item[D)] eine mechanische Kraft
    \end{itemize}

\ifshowsolutions
Lösung: A
\fi

  \item Warum kann sich die Spannung an einem idealen Kondensator nicht sprunghaft ändern?
    \begin{itemize}[label=$\square$]
      \item[A)] Da der Kondensator Gleichstrom sperrt, sind schnelle Spannungsänderungen nicht möglich.
      \item[B)] Der ideale Kondensator verhindert grundsätzlich jede Änderung der gespeicherten Energie.
      \item[C)] Eine sprunghafte Spannungsänderung würde einen unendlich großen Strom erfordern.
      \item[D)] Die Spannung an einem Kondensator ist durch die Kapazität nach oben begrenzt.
    \end{itemize}

\ifshowsolutions
Lösung: C
\fi

  \item Bei konstanter Gleichspannung verhält sich eine ideale Spule im stationären Zustand wie:
    \begin{itemize}[label=$\square$]
      \item[A)] Stromunterbrechung
      \item[B)] Stromquelle
      \item[C)] Spannungsquelle
      \item[D)] Kurzschluss
    \end{itemize}

\ifshowsolutions
Lösung: D
\fi

  \item Ein Kondensator wird über einen ohmschen Widerstand entladen. Welche Aussage beschreibt den Entladevorgang korrekt?
    \begin{itemize}[label=$\square$]
      \item[A)] Die Spannung am Kondensator fällt exponentiell ab, der Strom nimmt exponentiell zu.
      \item[B)] Die Spannung am Kondensator nimmt exponentiell zu, der Strom fällt exponentiell ab.
      \item[C)] Die Spannung am Kondensator fällt exponentiell ab, der Strom fällt exponentiell ab.
      \item[D)] Die Spannung am Kondensator nimmt exponentiell zu, der Strom nimmt exponentiell zu.
    \end{itemize}

\ifshowsolutions
Lösung: C
\fi

  \item Warum wird bei Netzspannungen üblicherweise der Effektivwert angegeben?
    \begin{itemize}[label=$\square$]
      \item[A)] Weil er kleiner ist als der Spitzenwert.
      \item[B)] Weil er die Leistungswirkung der Spannung in ohmschen Lasten beschreibt.
      \item[C)] Weil er unabhängig von der Frequenz ist.
      \item[D)] Weil er den zeitlichen Verlauf der Spannung vollständig beschreibt.
    \end{itemize}

\ifshowsolutions
Lösung: B
\fi

  \item Bei einer idealen Spule im Wechselstromkreis gilt:
    \begin{itemize}[label=$\square$]
      \item[A)] Der Strom eilt der Spannung voraus
      \item[B)] Der Strom hinkt der Spannung hinterher
      \item[C)] Strom und Spannung sind um 45° gegeneinander verschoben
      \item[D)] Strom und Spannung sind gleichphasig
    \end{itemize}

\ifshowsolutions
Lösung: B
\fi

  \item Welche physikalische Ursache führt dazu, dass bei einem Kondensator der Strom der Spannung vorausläuft?
    \begin{itemize}[label=$\square$]
      \item[A)] Die Spannung bewirkt einen sofortigen Stromfluss durch das Dielektrikum.
      \item[B)] Die Spannung wird durch ohmsche Verluste im Kondensator verzögert.
      \item[C)] Magnetische Effekte bestimmen das Verhalten des Kondensators.
      \item[D)] Der Strom lädt den Kondensator, wodurch sich die Spannung verzögert aufbaut.
    \end{itemize}

\ifshowsolutions
Lösung: D
\fi

  \item Welche physikalische Ursache führt dazu, dass in einem Wechselstromkreis Blindleistung auftritt?
    \begin{itemize}[label=$\square$]
      \item[A)] Elektrische Energie wird periodisch in Feldern gespeichert und wieder abgegeben.
      \item[B)] Elektrische Energie wird periodisch in Wärme gespeichert und wieder abgegeben.
      \item[C)] Elektrische Energie wird periodisch an die Quelle zurückgeführt und gespeichert.
      \item[D)] Elektrische Energie wird periodisch im Leiter transportiert und verzögert.
    \end{itemize}

\ifshowsolutions
Lösung: A
\fi

  \item Welche Filtercharakteristik zeigt eine Parallelschaltung aus Spule und Kondensator, wenn sie in Serie zu einer Last geschaltet wird?
    \begin{itemize}[label=$\square$]
      \item[A)] Tiefpass
      \item[B)] Hochpass
      \item[C)] Bandpass
      \item[D)] Bandsperre
    \end{itemize}

\ifshowsolutions
Lösung: D
\fi

  \item Ein symmetrisches Drehstromsystem:
    \begin{itemize}[label=$\square$]
      \item[A)] erzeugt ein rotierendes elektrisches Feld
      \item[B)] benötigt nur bei unsymmetrischer Last einen Neutralleiter
      \item[C)] arbeitet mit Gleichstrom
      \item[D)] hat keine Spannung zwischen den Außenleitern
    \end{itemize}

\ifshowsolutions
Lösung: B
\fi

\end{enumerate}

\end{document}